Глава отвечает на вопрос, оставшийся с главы~6: почему квазиньютоновские методы сохраняют сверхлинейную сходимость метода Ньютона, хотя $B_k$ не обязана сходиться к гессиану. Ответ --- \hyperref[note:D65FB534-B7F1-4194-8E4B-D4ACAA89E6D2]{критерий Денниса--Море}: достаточно приближать гессиан вдоль шага $d_k$; вместе с условием секущей это даёт \hyperref[note:A9E5C7FA-46FE-460E-A1B9-5B7CD72DB5B4]{сверхлинейную сходимость}. Вторая тема --- функция-сумма $F = \frac{1}{N}\sum f_i$ и \hyperref[note:A4769F74-2E6F-4CA3-89E4-7B01881D8109]{стохастический градиентный спуск}, который на каждом шаге берёт градиент одного слагаемого; для выпуклой $F$ \hyperref[note:58D1735C-5A70-4DA0-B522-B4595CB4494C]{оценка скорости} верна в среднем и для усреднённой точки $\overline{x}_k$.
